\section{Introduction}

A group is \emph{polycyclic} if it has a finite subnormal series with cyclic factors. It is \emph{virtually polycyclic} if it contains a polycyclic subgroup of finite index. We prove that virtual polycyclicity can be recognized from the large-scale geometry of a word metric.

A $(K,C)$-\emph{quasi-isometry} $F:X\to Y$, with $K\geq1$ and $C\geq0$, satisfies
\[
 K^{-1}d_X(x,x')-C\leq d_Y(F(x),F(x'))\leq Kd_X(x,x')+C
\]
for all $x,x'\in X$, and every point of $Y$ is within distance $C$ of $F(X)$. Different finite generating sets give quasi-isometric word metrics. 
\begin{theorem}\label{thm:recognition}
Let $P$ be a finitely generated virtually polycyclic group, and let $H$ be any finitely generated group. If $H$ and $P$ are quasi-isometric, then $H$ is virtually polycyclic.
\end{theorem}

The group $H$ in Theorem~\ref{thm:recognition} is subject to no algebraic hypothesis beyond finite generation.

The virtual lattice characterization of polycyclic groups gives the following equivalent formulation. A lattice is a discrete subgroup of finite covolume; it is uniform if the quotient is compact.

\begin{corollary}\label{cor:lattice-recognition}
If a finitely generated group $H$ is quasi-isometric to a lattice in a connected simply connected solvable Lie group, then a finite-index subgroup of $H$ is a uniform lattice in some connected simply connected solvable Lie group.
\end{corollary}

The ambient Lie group in this conclusion may change. The precise lattice statements and finite-index reductions are given in the proof of Corollary~\ref{cor:lattice-recognition}.

\subsection{The recognition problem}

Eskin, Fisher and Whyte formulated the lattice-recognition conjecture
and its equivalent polycyclic formulation in
\cite[Conjecture~1.2 and Remark~(3), p.~929]{EskinFisherWhyte2007}.
Gromov's theorem on groups of polynomial growth settles the virtually
nilpotent case \cite{Gromov1981}. For the full class, Shalom proved that
a finitely generated group quasi-isometric to an infinite polycyclic
group has a finite-index subgroup with infinite abelianization
\cite[Theorem~1.3]{Shalom2004}. The remaining recognition problem asks
for the algebraic structure of an arbitrary finitely generated
comparison group, on which solvability is not assumed.

Height rigidity has been central to the exponential-growth case.
Farb and Mosher formulated horizontal-preservation conjectures for
abelian-by-cyclic models and the recognition conjecture for
$\mathrm{Sol}$ \cite[Section~5]{FarbMosher2000}.
Eskin, Fisher and Whyte established recognition for $\mathrm{Sol}$
by coarse differentiation \cite[Theorem~1.1]{EskinFisherWhyte2013}.
Peng extended this analysis to nondegenerate unimodular split
abelian-by-abelian Lie groups \cite{Peng2011I,Peng2011II}.
Her argument makes local height maps affine and compares them across
boxes; it gives a description of quasi-isometries as well as
virtual-polycyclic recognition
\cite[Theorem~5.3.6 and Corollary~5.3.9]{Peng2011II}.
This class already includes nontrivial Jordan blocks.

For several of these models, height control induces actions on
boundaries, where conjugation theorems give stronger rigidity
conclusions.
Dymarz carried this out for diagonalizable determinant-one models
with no unit-modulus eigenvalues \cite{Dymarz2010}.
Dymarz, Fisher and Xie prove a conjugation theorem for suitable
nilpotent boundaries and apply it to SOL-like groups with possibly
nonabelian nilpotent factors
\cite[Theorems~1.2 and~1.5]{DymarzFisherXie2025}.
Their Section~9 repairs the conjugation step in
\cite[Section~3.3]{Dymarz2010} using amenability. Their recognition
result retains the structural and amenability hypotheses specified
in their Theorem~1.5.

Grayevsky and Pallier obtain rigidity of $\mathrm{Sol}_5$ and its
lattices \cite[Theorem~C]{GrayevskyPallier2025}.
A different strengthening of the geometric hypothesis is commability:
two locally compact groups are commable if they are joined by a finite
zigzag of continuous proper homomorphisms with cocompact image
\cite[Remark~4.C.13]{CornulierDeLaHarpe2016}.
Le Boudec proves virtual-polycyclic recognition under this hypothesis
\cite[Theorem~7 and Corollary~3.7]{LeBoudec2025}.

Theorem~\ref{thm:recognition} resolves recognition for the full class
of virtually polycyclic groups. We establish uniform height control
for the general solvable Lie models described below. Combined with
amenability, this control yields elementary amenability of the
comparison group; finiteness and duality then yield virtual polycyclicity.

\subsection{Uniform bounded height}

We use connected simply connected \emph{real-triangulable} Lie groups: the adjoint representation of their Lie algebra admits a simultaneously upper triangular form over $\R$. Fix a left invariant Riemannian metric on each such group. A group is \emph{unimodular} if its left Haar measure is also right invariant. The relevant quotient has a canonical definition.

\begin{definition}\label{def:canonical-height}
Let $\mathfrak g$ be the Lie algebra of a real-triangulable group $G$. Denote the diagonal characters of a simultaneous real triangular form of $\ad$ by $\chi_1,\ldots,\chi_s$, including repetitions and zero characters. Put
\[
 \mathfrak k=\bigcap_{j=1}^s\ker\chi_j,
 \qquad W=\mathfrak g/\mathfrak k.
\]
The \emph{exponential height map} is the homomorphism $\pi:G\to(W,+)$ whose differential is the quotient map.
\end{definition}

For the zero-dimensional algebra, the empty intersection is $\mathfrak g$. The characters vanish on $[\mathfrak g,\mathfrak g]$, so the quotient is abelian. Their multiset is intrinsic to the adjoint representation; hence $\mathfrak k$ is independent of the triangular basis. Simple connectedness gives the homomorphism $\pi$. Fix any norm on $W$.

\begin{theorem}[Uniform bounded height]\label{thm:bounded-height}
Let $G$ be a connected simply connected real-triangulable unimodular Lie group, with a fixed left invariant Riemannian metric, exponential height map $\pi:G\to W$, and a fixed norm on $W$. There is a finite subgroup $\A_G\leq\GL(W)$ such that, for every $K\geq1$ and $C\geq0$, there is a constant $B=B(G,K,C)$ with the following property. Every $(K,C)$ self quasi-isometry $F$ of $G$ has an $A_F\in\A_G$ satisfying
\begin{equation}\label{eq:height-main}
 \bigl\|\pi F(g)-\pi F(1)-A_F\pi(g)\bigr\|\leq B
 \qquad(g\in G).
\end{equation}
\end{theorem}

The finite group $\A_G$ will be constructed from the exponential rates of bracket-generated volume in Lemma~\ref{lem:support-symmetries}. It depends on $G$, whereas the error bound depends additionally on the quasi-isometry constants. When $W=0$, the height statement is vacuous and $G$ is nilpotent; this case is treated separately in recognition.

\subsection{How height gives recognition}

We first explain how Theorem~\ref{thm:bounded-height} implies recognition; the complete deduction appears in Section~\ref{sec:recognition}. A virtually polycyclic group is quasi-isometric to a connected simply connected real-triangulable unimodular Lie model $G$. Transporting the left translation action of $H$ across a quasi-isometry gives maps of $G$ with uniform quasi-isometry bounds that respect multiplication up to a uniformly bounded error. This is a \emph{uniform quasi-action}. Equation~\eqref{eq:height-main} assigns a linear part in the finite group $\A_G$ to each element of $H$. After passing to a finite-index subgroup $H_0$, all linear parts are the identity. The group $H$ inherits amenability from its virtually polycyclic comparison group by quasi-isometry. Amenability then corrects the bounded defect in the height translations to a homomorphism $\tau:H_0\to W$.

Every finitely generated subgroup of $\ker\tau$ has polynomial growth: its word paths have bounded height, and the volume of possible endpoints grows polynomially. Gromov's theorem makes each such subgroup virtually nilpotent. Thus $\ker\tau$, and consequently $H$, is elementary amenable. Finiteness properties and Poincar\'e duality inherited from the comparison group, together with the structure of elementary amenable groups, then imply that a finite-index subgroup of $H$ is polycyclic. The bounded error is essential: a sublinear height bound alone would not give the bounded-height word paths used to prove polynomial growth in the kernel.

\subsection{The height argument}

The geometric meaning of height is already visible in
$\mathrm{Sol}=\R^2\rtimes\R$, where $t\in\R$ acts on the two coordinate
directions by $e^t$ and $e^{-t}$.  A bounded horizontal move at height $t$
can change the first coordinate by order $e^t$ and the second by order
$e^{-t}$.  To reach many endpoints, a path can make its first-coordinate
move high up and its second-coordinate move low down.  If its allowed
heights form $r[a,b]$, the resulting rectangle has area of order
$e^{r(b-a)}$.  Thus the height range determines an exponential rate of
accessible volume.

For a general Lie model, the corresponding rate is a function $P(E)$ of
the compact convex allowed height set $E\subset W$.  Its definition in
Section~\ref{sec:boxes} includes every bracket direction generated by the
expanding directions.  Paths of length $O(r)$ from one starting point,
visiting the appropriate support positions and returning to a common
height, reach
$e^{rP(E)+o(r)}$ endpoints separated at a fixed scale.  Their images also have length
$O(r)$.  If those images remain at heights in $rE'$ up to sublinear
errors, the image paths can
reach at most $e^{rP(E')+o(r)}$ separated endpoints.  A quasi-isometry
therefore forces $P(E)\le P(E')$.  Applying this comparison in both
directions will restrict an affine height slope to the finite group
$\mathcal A_G$.

The proof must first produce such slopes and then make them independent
of the averaging, chart, and scale used to find them.  Choose the
nilpotent groups $N,D$ of Section~\ref{sec:models}, with $N$ normal and
$G=ND$.  The expression $g=nd$ need not be unique.  We study the height
function on the sets $xnD$: for fixed $x\in G$ and $n\in N$, rescale
$d\mapsto\pi F(xnd)$ on larger and larger intrinsic $D$-balls.  These are
the $D$-charts, or sheets indexed by $n$.  Endpoint counts instead use
the slices $xNd$ at a fixed $D$-position, where the parameterization by
$N$ is injective.

The geometric proof has three stages. First we study a compact space of
normalized quasi-isometries: each map is translated so that it sends the
identity to the identity. Moving the basepoint and recentering defines
an action on this space. This construction follows Shalom's
measured-coupling viewpoint \cite[Sections~2.1--2.2]{Shalom2004}.
For each ergodic invariant probability measure, the reduced-cohomology
vanishing theorem of Cornulier and Tessera
\cite[Corollary~1.9]{CornulierTessera2020} supplies a linear height slope.
The volume comparison above restricts this slope to $\mathcal A_G$ once
it is applied to the map and its inverse. For the reverse comparison we
need the inverse of the same slope; a coupling constructed from Haar
measure provides it.

Second we pass from invariant measures to arbitrary maps. We sample
large $D$-charts and rescale their height functions. Stationary
increments constrain each limiting profile, while accessible endpoint
volume prevents it from folding. Together these facts make every
supported limiting profile affine on the chart interior. The volume
argument still works after restricting the switch parameters to a set
retaining a subexponential fraction of their measure: a polynomial-image estimate retains the
required exponential endpoint volume.
Different sheets could nevertheless have different affine maps.
If neighboring profiles had different linear parts, a limiting
comparison at supporting heights would force a suitable source height
segment to have opposite target directions on the two sheets.
Alternating moves on the sheets would then produce a target volume
exponent twice the source exponent, contradicting the volume comparison.
Haar multipliers in suitable nilpotent
quotients then detect, and eliminate, differences of their translations.
Comparison across scales and basepoints gives one global slope with
a uniform sublinear error.

Finally we improve this error to a bounded one. Contracting rays define
maps between nilpotent quotient boundaries. The Haar volumes of their
images record height; averaging translated sets then bounds its error.
Boundary conjugation theorems, such as those in
\cite{Dymarz2010,DymarzFisherXie2025}, produce rigid actions from height
control. Here, starting from the uniform sublinear estimate, the
boundary maps bound the remaining height error through their image
measures. The recognition argument can then use
the bounded-height kernel and the finiteness and duality theorems.

The progression from weak height control to bounded error, followed by
an application to a transported group action, already occurs in
\cite[Sections~6.1--6.2 and Theorem~7.3]{EskinFisherWhyte2012}.
The present proof establishes the estimate for the full class in
Theorem~\ref{thm:bounded-height}. Its volume calculations retain both
the bracket directions in $N$ and the multiplicity of the surjection
$N\rtimes D\to G$, so no complementary height subgroup is required.

Sections~\ref{sec:models}--\ref{sec:polynomial-volume} develop the Lie
coordinates, volume estimates, and restricted-product lemma.
Section~\ref{sec:stationary} identifies stationary slopes. The profile
and neighboring-chart arguments lead to the uniform sublinear estimate
in Section~\ref{sec:uniformity}; Section~\ref{sec:boundary} proves bounded
height. Section~\ref{sec:recognition} completes the group-theoretic
deduction.
